\documentclass{article}

% if you need to pass options to natbib, use, e.g.:
%     \PassOptionsToPackage{numbers, compress}{natbib}
% before loading neurips_2026

% The authors should use one of these tracks.
% Before accepting by the NeurIPS conference, select one of the options below.
% 0. "default" for submission
\usepackage{neurips_2026}

\usepackage[utf8]{inputenc} % allow utf-8 input
\usepackage[T1]{fontenc}    % use 8-bit T1 fonts
\usepackage{hyperref}       % hyperlinks
\usepackage{url}            % simple URL typesetting
\usepackage{booktabs}       % professional-quality tables
\usepackage{amsfonts}       % blackboard math symbols
\usepackage{nicefrac}       % compact symbols for 1/2, etc.
\usepackage{microtype}      % microtypography
\usepackage{xcolor}         % colors

\title{FinRL-DeepSeek Risk-First: Structurally Constraining RL Agents with LLM Semantic Signals}


\author{%
  Seth Ndinga-Ndinga \\
  Aivancity School for Technology, Business \& Society \\
  Paris-Cachan, France \\
  \texttt{seth.ndinga-ndinga@aivancity.ai} \\
}


\begin{document}


\maketitle


\begin{abstract}
Portfolio optimization via deep reinforcement learning (DRL) consistently fails during market crashes because it focuses exclusively on return maximization. Large language models (LLMs) can read financial news to anticipate these crashes, but our experiments show that a standard DRL agent ignores semantic signals in favor of price momentum. We present \textit{Risk-First}, an architecture that mathematically forces the agent to follow those alerts. The system has three components: a variance filter to discard LLM hallucinations, a reward penalty for dangerous exposure (\textit{Reward Shaping}), and a deterministic circuit breaker that forces asset liquidation. Tested on the NASDAQ market using the FNSPID dataset and DeepSeek-V3, these constraints reduce tail risk (Max Drawdown) from -58.37\% to -56.09\% while increasing returns, showing that the DRL architecture must be structurally constrained to make LLM predictions useful.
\end{abstract}

\section{Introduction}

Deep reinforcement learning (DRL) handles stock trading well, but standard algorithms like PPO maximize returns without managing loss severity. During market crashes, this risk-neutral behavior produces capital drawdowns that are not acceptable in production.

Models like DeepSeek-V3 can extract real-time sentiment and risk indicators from financial news. The FinRL-DeepSeek architecture \cite{benhenda2025} feeds these signals into the agent's state. The problem: giving risk information to a risk-neutral agent does not produce cautious behavior. When price momentum is strong, the neural network ignores the textual alert and chases short-term gains.

To force the agent to follow these alerts, the \textit{Risk-First} architecture modifies the learning process through three modules: (1) An empirical confidence filter that neutralizes ambiguous LLM predictions; (2) A reward penalty (\textit{Reward Shaping}) that mathematically reduces gains when the agent holds a position in an asset the LLM flags as risky; and (3) A mechanical circuit breaker that forces liquidation and blocks buys when a danger threshold is crossed. The ablation study shows that these constraints cut crash exposure and outperform standard DRL approaches.

\section{Related Work}

\subsection{Deep Reinforcement Learning in Quantitative Trading}

Quantitative trading has historically relied on heuristics, moving average crossovers, and supervised learning models. DRL addresses this by framing trading as a Markov Decision Process (MDP). The open-source FinRL library \cite{liu2020} standardized DRL agent development in finance. Proximal Policy Optimization (PPO) \cite{schulman2017} remains the reference algorithm. While PPO performs well in bull markets, it ignores loss severity during crashes. Similar RL approaches have been successfully applied to other financial domains, such as automated liquidity provisioning in decentralized finance \cite{xu2025}.

\subsection{Large Language Models for Financial Forecasting}

LLMs now allow finer-grained text analysis, extracting sentiment and risk indicators in a zero-shot setting on datasets like FNSPID \cite{dong2024}. However, financial LLMs are susceptible to look-ahead bias, requiring rigorous evaluation against benchmarks like Look-Ahead-Bench \cite{benhenda2026lookaheadbench}. Furthermore, the predictive capabilities of LLMs in finance are being mapped through benchmarks like VCBench \cite{chen2025vcbenchbenchmarkingllmsventure} for venture capital, YCBench \cite{benhenda2026ycbench} for startup forecasting, and FutureX \cite{zeng2025, chandak2026} for open-ended prediction. Benhenda \cite{benhenda2025} introduced the FinRL-DeepSeek architecture, but the original implementation consumes scores at face value, ignoring predictive uncertainty.

\subsection{Risk-Sensitive Reinforcement Learning (CVaR)}

Risk-sensitive RL addresses the weaknesses of risk-neutral algorithms. While finance commonly uses Value at Risk (VaR), Conditional Value at Risk (CVaR) \cite{rockafellar2000} is a stricter measure quantifying tail risk. The CPPO architecture embeds CVaR into Policy Gradient algorithms. The hybrid approach in this paper combines CPPO's retrospective statistical rigor with the LLM's forward-looking semantic signals.

\section{Methodology}

\subsection{Background: Markov Decision Process (MDP)}

The portfolio allocation problem is modeled as an MDP $(\mathcal{S}, \mathcal{A}, \mathcal{P}, \mathcal{R}, \gamma)$. At each time step $t$, the state $s_t$ contains historical asset data and LLM scores. The action $a_t$ is a continuous vector defining asset weights. The reward $r_t$ is the change in total portfolio value. The agent finds a policy $\pi_\theta(a_t|s_t)$ that maximizes cumulative expected rewards.

\subsection{Module 1: Semantic Confidence Filter (LLM Confidence Gate)}

The first module implements an empirical confidence filter. If the confidence level (derived from multiple query variance) is below $\tau = 0.7$, the signal is neutralised to 3 (on a 1 to 5 scale). This filter prevents overreaction to contradictory financial news.

\subsection{Module 2: Reward Penalty (Risk-Adjusted Reward Shaping)}

The second module modifies the reward function. A penalty proportional to the semantic risk score and the asset's weight is subtracted from the raw reward. This forces the agent to trade off return against immediate semantic risk.

\subsection{Module 3: Mechanical Circuit Breaker (Emergency Circuit Breaker)}

The third module acts as a deterministic circuit breaker. When semantic signals reach a critical threshold (risk $\ge 4$ and sentiment $\le 2$), the system cancels buy orders and mechanically reduces existing positions. This protects capital during extreme events.

\section{Experiments (Ablation Study)}

\subsection{Experimental Setup}

The environment was built on the Gymnasium API. Market data covers major NASDAQ assets. Text signals were extracted from the FNSPID dataset. DeepSeek-V3 was applied zero-shot. Neural networks were trained for 100 epochs, with chronological train/test splits.

\subsection{Evaluation Metrics}

Performance was measured using Cumulative Return (\%), Max Drawdown (\%), Sharpe Ratio, and Rachev Ratio ($\alpha = 0.05$).

\subsection{Ablation Configurations}

Six configurations were tested incrementally:
\begin{itemize}
    \item \textbf{A (PPO Baseline)}: Standard PPO.
    \item \textbf{B (PPO + LLM)}: PPO receiving raw LLM scores.
    \item \textbf{C (CPPO Baseline)}: CPPO constrained by purely statistical CVaR.
    \item \textbf{D (CPPO + LLM Confidence Gate)}: CPPO with Module 1.
    \item \textbf{E (CPPO + Reward Shaping)}: Adds Module 2.
    \item \textbf{F (Full Risk-First Architecture)}: Full model, adds Module 3.
\end{itemize}

\subsection{Ablation Study Results}

\begin{table}[h]
  \caption{Comparative performance across ablation configurations.}
  \label{performance-table}
  \centering
  \begin{tabular}{lcccc}
    \toprule
    Configuration & Cum. Return & Sharpe Ratio & Rachev Ratio & Max Drawdown \\
    \midrule
    \textbf{A} (PPO Baseline)    & 160.08\% & 0.726 & 0.915 & -58.37\% \\
    \textbf{B} (PPO + LLM)       & 160.08\% & 0.726 & 0.915 & -58.37\% \\
    \textbf{C} (CPPO Baseline)   & 160.08\% & 0.726 & 0.915 & -58.37\% \\
    \textbf{D} (CPPO + LLM Gate)     & 168.71\% & 0.752 & 0.926 & -57.37\% \\
    \textbf{E} (CPPO + Shaping)  & 179.77\% & 0.773 & 0.939 & -56.51\% \\
    \textbf{F} (Risk-First Full) & 178.97\% & 0.784 & 0.928 & -56.09\% \\
    \bottomrule
  \end{tabular}
\end{table}

The table shows that the LLM only helps when the agent is structurally forced to act on it:
\begin{itemize}
    \item \textbf{Failure of standard approaches (Configs A, B, C):} The baseline models do not protect against crashes. Config B shows that giving raw LLM text to the agent changes nothing: the network ignores the semantic signal and follows price momentum. Config C shows that the statistical CVaR constraint stays inactive because historical data alone cannot anticipate the crash.
    \item \textbf{Impact of the semantic filter and reward penalty (Configs D and E):} Activating the confidence filter removes textual noise and raises returns to 168.71\%. Adding Reward Shaping mathematically reduces gains when the agent ignores a high-risk alert, pushing return to 179.77\% and cutting Max Drawdown to -56.51\%.
    \item \textbf{Effect of the mechanical circuit breaker (Config F):} Adding the circuit breaker produces the highest Sharpe ratio (0.784) and the lowest Max Drawdown (-56.09\%). Cumulative return falls slightly compared to Config E (178.97\% vs 179.77\%): the circuit breaker blocks buys and trims existing positions during extreme-risk days, compressing raw performance but improving the risk-adjusted profile.
\end{itemize}

\section{Conclusion}

The ablation study shows that the \textit{Risk-First} architecture (Config F) produces concrete improvements in risk management during bear market conditions. By combining semantic confidence filtering, reward modulation, and the deterministic circuit breaker, the full model reaches the highest Sharpe ratio (0.784) and the lowest Max Drawdown (-56.09\%) across all tested configurations, outperforming both PPO and CPPO baselines. The small drop in cumulative return compared to Config E reflects the circuit breaker acting defensively, trading a marginal share of raw gain for better downside protection. Future work should run multi-seed evaluations to confirm statistical significance and test the architecture on more volatile asset classes such as cryptocurrencies.

\bibliographystyle{plain}
\bibliography{references}

\end{document}
